\section{Limitations} \label{sec:limitations}

Verus currently only supports mutable borrows (\verb`&mut`) of data passed as arguments to a function call:
mutable references in return values and explicit borrows on the right-hand side of assignments are not
supported.
We believe that adding more complete support with an approach similar to Creusot's is mainly a
matter of syntax, interface design, and engineering.

Unlike tools that re-encode ownership properties (e.g., with separation logic in Viper~\cite{DBLP:conf/vmcai/0001SS16}),
Verus relies on borrow-checking rules and hence cannot reason about traditional Rust unsafe code.
This may limit its applicability in applications that heavily rely on unsafe,
e.g., when direct memory manipulation is required to communicate with memory-mapped devices.
\autoref{sec:unsafe}, \autoref{sec:linear-ghost}, and \autoref{sec:dlist} discuss
encapsulations that alleviate the need for unsafe in certain contexts.

Verus is closely tied to Rust's type system,
which is more limited in some ways than the dependent type systems of Coq and F*.
This may preclude some more sophisticated styles of structuring proofs that are supported by Coq and F*.
While the limitation on mutable borrows will be lifted in the near future, limitations tied to Rust's type
system are imposed by Verus's design choices.

